% 出席確認クイズ 第13回　AIのリスクとガバナンス
%   attendance/attendance13.html から生成（解説は本ガイド用に加筆）。

\quizq{1}{AIの「バイアス」の主な原因として適切なものはどれでしょうか?}%
  {電力不足}%
  {画面の大きさ}%
  {\correct{学習データの偏り}}%
  {通信速度}%
  {正解は (c)。学習データに偏りがあると、AIの判断もその偏りを引き継ぎます。たとえば過去の採用データに性別の偏りがあれば、そこから学習した選考モデルも同じ偏りを持ちます。(a)(b)(d) はバイアスの原因ではありません。}

\quizq{2}{「シャドーAI」の説明として適切なものはどれでしょうか?}%
  {公式に承認されたAI}%
  {\correct{組織の許可なく従業員がAIツールを業務に使うこと}}%
  {夜間に稼働するAI}%
  {画像の影を処理するAI}%
  {正解は (b)。シャドーAIは、組織が把握・許可していない AI ツールを従業員が業務に使うことで、機密情報の漏えいなどのリスクがあります。シャドーIT（無許可の IT 利用）の派生概念です。禁止するだけでは地下化するため、安全な公式ツールの提供と教育が対策になります。}

\quizq{3}{EUのAI規制(AI Act)の基本的な考え方として適切なものはどれでしょうか?}%
  {規制は一切しない}%
  {企業の規模だけで決める}%
  {すべてのAIを禁止する}%
  {\correct{リスクの高さに応じて義務や規制の強さを変える}}%
  {正解は (d)。EU AI Act は、AIの用途を「許容できないリスク」「高リスク」「限定的リスク」「最小リスク」に分類し、リスクの高さに応じた義務を課すリスクベースの規制です。(a)(b)(c) はいずれも実際の枠組みとは異なります。}

\quizq{4}{生成AIと著作権に関する説明として適切なものはどれでしょうか?}%
  {\correct{出力が既存の著作物に類似していると、著作権侵害になる可能性がある}}%
  {著作権はAIには関係ない}%
  {著作権はAIが持つ}%
  {AIが作ったものはすべて自由に使える}%
  {正解は (a)。生成AIの出力であっても、既存の著作物と類似し、それに依拠していると判断されれば著作権侵害になる可能性があります。(b)(d) は誤りで、(c) 現行法では AI は著作者になれません。学習段階と利用（出力）段階で考え方が異なる点にも注意が必要です。}

\quizq{5}{企業のAIガバナンスに含まれるものとして適切なものはどれでしょうか?}%
  {ルールを作らない}%
  {AIの利用を隠す}%
  {全社員に無制限の利用を許可する}%
  {\correct{利用方針の策定、リスク評価、責任者・体制の明確化}}%
  {正解は (d)。AIガバナンスは、利用方針の策定、リスク評価、責任者・体制の明確化、教育、監査などを整え、AIを適切に活用するためのしくみです。(a)(c) はガバナンスの不在、(b) は透明性の否定です。}
